\documentclass[11pt,a4paper]{article}
\usepackage[margin=1in]{geometry}
\usepackage[T1]{fontenc}
\usepackage[utf8]{inputenc}
\usepackage{lmodern}
\usepackage{xcolor}
\usepackage{hyperref}
\usepackage{enumitem}
\usepackage{listings}
\usepackage{titlesec}
\usepackage{longtable}

\hypersetup{
  colorlinks=true,
  linkcolor=blue,
  urlcolor=blue,
  pdftitle={Qlanka-Pro Interview Q and A},
  pdfauthor={Qlanka-Pro Candidate Notes}
}

\lstdefinestyle{code}{
  basicstyle=\ttfamily\small,
  breaklines=true,
  frame=single,
  columns=fullflexible,
  keepspaces=true,
  showstringspaces=false,
  keywordstyle=\color{blue!70!black},
  commentstyle=\color{green!40!black},
  stringstyle=\color{orange!60!black}
}

\title{Qlanka-Pro Full Interview Q and A\\\large WSO2, Microservices, Security, and Operations}
\author{Project-Based Preparation Guide}
\date{April 2026}

\begin{document}
\maketitle
\tableofcontents
\newpage

\section{How to Use This Document}
\begin{itemize}[leftmargin=*]
  \item This is a project-grounded interview script for Qlanka-Pro.
  \item Answers are written in a practical style suitable for technical interviews.
  \item Code snippets are adapted from the real repository and can be discussed line-by-line.
  \item For each question, start with a 20--30 second summary, then go deep only if asked.
\end{itemize}

\section{Project Snapshot (30-Second Pitch)}
\textbf{Q: Give me a quick summary of Qlanka-Pro.}\\
\textbf{A:} Qlanka-Pro is a smart queue platform built as .NET 8 microservices with a React frontend. It uses a gateway pattern, separate domain services (Identity, Queue, ServiceCenter, Notification), and Prometheus/Grafana for observability. We integrated WSO2 Identity concepts into authentication paths using OAuth2 token issuance and SCIM-style provisioning flows, while adding dynamic token validation to support both legacy local JWT and WSO2-issued tokens during migration.

\section{WSO2 and Identity Security (Most Critical)}

\subsection{Walk me through your WSO2 integration}
\textbf{Q: How did you integrate WSO2 Identity Server into Qlanka-Pro? What were the main challenges?}\\
\textbf{A:}
\begin{itemize}[leftmargin=*]
  \item In Identity service, login delegates credential validation and token issuance to WSO2 using OAuth2 password grant when \texttt{Wso2:Enabled=true}.
  \item Registration triggers provisioning to WSO2 via SCIM endpoint \texttt{/scim2/Users}.
  \item Gateway and microservices use dynamic JWT schemes: local tokens for backward compatibility and WSO2 authority-based validation for migrated flows.
  \item Main challenges:
  \begin{itemize}[leftmargin=*]
    \item Zero-downtime migration from local JWT to external IdP.
    \item Claim mapping differences (role claim naming and issuer/audience alignment).
    \item Environment-specific TLS and metadata issues in development.
    \item Maintaining consistency between app DB users and WSO2 provisioned users.
  \end{itemize}
\end{itemize}

\textbf{Important snippet: Identity login delegates to WSO2}\newline
\begin{lstlisting}[style=code,language=C]
Wso2TokenResult? wso2Token = null;
if (IsWso2Enabled())
{
    // Delegate credential validation and token issuance to WSO2 when enabled.
    wso2Token = await _wso2Identity.RequestTokenAsync(dto.Username, dto.Password);
}
else
{
    // Local fallback preserves compatibility for environments without WSO2.
    if (!BCrypt.Net.BCrypt.Verify(dto.Password, user.PasswordHash))
        throw new InvalidCredentialsException();
}

return new LoginResponseDto
{
    Token = wso2Token?.AccessToken ?? GenerateJwt(user),
    RefreshToken = wso2Token?.RefreshToken ?? GenerateRefreshToken(),
    ExpiresIn = wso2Token?.ExpiresIn ?? expiryMins * 60,
    Role = user.Role.ToLowerInvariant(),
    CounterId = counterId
};
\end{lstlisting}

\subsection{Why ROPC over Authorization Code?}
\textbf{Q: Why did you choose ROPC over Authorization Code flow? What are trade-offs?}\\
\textbf{A:}
\begin{itemize}[leftmargin=*]
  \item We used ROPC as a migration bridge because the existing UX was backend-driven username/password login.
  \item It reduced integration complexity and allowed fast rollout without immediate frontend redirect/callback redesign.
  \item Trade-offs:
  \begin{itemize}[leftmargin=*]
    \item Weaker security model than Authorization Code + PKCE for browser/mobile apps.
    \item Credentials pass through application tier.
    \item Less future-proof for federated SSO patterns.
  \end{itemize}
  \item Strategic direction: move to Authorization Code + PKCE for long-term best practice.
\end{itemize}

\textbf{Important snippet: ROPC request payload}\newline
\begin{lstlisting}[style=code,language=C]
using var body = new FormUrlEncodedContent(new Dictionary<string, string>
{
    ["grant_type"] = "password",
    ["username"] = username,
    ["password"] = password,
    ["scope"] = "openid profile",
    ["client_id"] = clientId,
    ["client_secret"] = clientSecret
});
\end{lstlisting}

\subsection{JWKS-based validation}
\textbf{Q: How does JWKS-based validation work and why is it better than symmetric key validation?}\\
\textbf{A:}
\begin{itemize}[leftmargin=*]
  \item With WSO2 mode enabled, services set \texttt{Authority} and validate tokens using IdP metadata/JWKS.
  \item The IdP signs tokens with private key; services validate with public keys.
  \item Advantages over symmetric key sharing:
  \begin{itemize}[leftmargin=*]
    \item No shared secret copied across microservices.
    \item Better key rotation behavior.
    \item Clear trust boundary: issuer signs, services verify.
  \end{itemize}
\end{itemize}

\textbf{Important snippet: WSO2 authority-based JWT validation}\newline
\begin{lstlisting}[style=code,language=C]
options.Authority = wso2Authority;
options.RequireHttpsMetadata = wso2RequireHttpsMetadata;
options.MapInboundClaims = false;

options.TokenValidationParameters = new TokenValidationParameters
{
    ValidateIssuer = true,
    ValidateAudience = true,
    ValidateLifetime = true,
    ValidateIssuerSigningKey = true,
    ValidIssuer = string.IsNullOrWhiteSpace(wso2ValidIssuer) ? null : wso2ValidIssuer,
    ValidAudience = wso2Audience,
    RoleClaimType = "role",
    ClockSkew = TimeSpan.Zero
};
\end{lstlisting}

\subsection{Token expiry and refresh}
\textbf{Q: How did you handle token expiry and refresh in your setup?}\\
\textbf{A:}
\begin{itemize}[leftmargin=*]
  \item Identity service returns access token, refresh token, and expiry metadata.
  \item Microservices are stateless and validate token on each request.
  \item Refresh logic is centralized at identity boundary, not inside domain services.
  \item This keeps services simple and reduces cross-service auth coupling.
\end{itemize}

\subsection{SCIM 2.0 provisioning}
\textbf{Q: What is SCIM 2.0 and how did you use it?}\\
\textbf{A:}
\begin{itemize}[leftmargin=*]
  \item SCIM 2.0 is a standard for user/group provisioning across systems.
  \item On registration, Qlanka-Pro calls WSO2 \texttt{/scim2/Users} with core attributes, role grouping, and officer center extension data.
  \item This reduces manual user management and keeps identity data synchronized.
\end{itemize}

\textbf{Important snippet: SCIM payload and request}\newline
\begin{lstlisting}[style=code,language=C]
var scimPayload = new Dictionary<string, object?>
{
    ["schemas"] = new[] { "urn:ietf:params:scim:schemas:core:2.0:User" },
    ["userName"] = username,
    ["password"] = password,
    ["emails"] = new[] { new { value = email, primary = true } },
    ["groups"] = new[] { new { display = role, value = role } }
};

if (centerId.HasValue)
{
    scimPayload["urn:scim:wso2:qlanka:1.0"] = new { centerId = centerId.Value.ToString() };
}
\end{lstlisting}

\subsection{API Manager vs reverse proxy}
\textbf{Q: How is WSO2 API Manager different from a simple reverse proxy?}\\
\textbf{A:}
\begin{itemize}[leftmargin=*]
  \item Reverse proxy (like YARP) mainly handles routing and pass-through concerns.
  \item API Manager adds governance: subscriptions, plans, consumer throttling, lifecycle/versioning, developer portal, and centralized policy control.
  \item In Qlanka-Pro, API publication metadata is prepared in \texttt{officer-api.yaml} with WSO2 extensions.
\end{itemize}

\subsection{Rate limiting in API Manager}
\textbf{Q: How did you configure rate limiting in WSO2 API Manager?}\\
\textbf{A:}
\begin{itemize}[leftmargin=*]
  \item We define per-operation throttling tiers in API definition metadata.
  \item Example: queue polling endpoints use a specific tier while command endpoints can use different policies.
  \item We also keep gateway-side limiter as defense-in-depth.
\end{itemize}

\textbf{Important snippet: WSO2 APIM throttling tiers}\newline
\begin{lstlisting}[style=code,language=]
'/api/counters/{counterId}/dashboard':
  get:
    x-throttling-tier: QueuePollingPolicy

'/api/counters/{counterId}/call-next':
  post:
    x-throttling-tier: Unlimited
\end{lstlisting}

\subsection{API lifecycle management}
\textbf{Q: What is API lifecycle management and how did you implement it?}\\
\textbf{A:}
\begin{itemize}[leftmargin=*]
  \item Lifecycle management controls stages such as create, publish, deprecate, and retire.
  \item We use versioned OpenAPI definitions and WSO2 extensions to support governed publication and evolution.
  \item This prevents uncontrolled breaking changes and improves consumer communication.
\end{itemize}

\section{Microservices and Architecture}

\subsection{Architecture layering}
\textbf{Q: Describe your Client -> API Manager -> Identity Server -> Microservices -> Database architecture. Why this layering?}\\
\textbf{A:}
\begin{itemize}[leftmargin=*]
  \item Client only talks to managed API surface.
  \item API entrypoint applies auth, throttling, routing policies.
  \item Identity centralizes authentication and token governance.
  \item Domain services isolate business capabilities and scale independently.
  \item Datastores are separated by domain to reduce coupling.
  \item This structure improves security, ownership boundaries, and deployability.
\end{itemize}

\textbf{Important snippet: docker-compose service boundaries}\newline
\begin{lstlisting}[style=code,language=]
services:
  gateway:
    depends_on: [identity, service-center, queue]
  identity:
    depends_on: [db-identity]
  service-center:
    depends_on: [db-service-center]
  queue:
    depends_on: [db-queue]
\end{lstlisting}

\subsection{REST vs gRPC}
\textbf{Q: How do your microservices communicate: REST or gRPC? When choose each?}\\
\textbf{A:}
\begin{itemize}[leftmargin=*]
  \item Current project communication is HTTP/REST based.
  \item REST fits API gateway integration, OpenAPI-driven governance, and broad interoperability.
  \item gRPC is preferred when low-latency high-throughput service-to-service communication or streaming is critical.
  \item Practical approach: REST for external/public paths, gRPC for internal high-performance paths if needed.
\end{itemize}

\subsection{Distributed tracing and observability}
\textbf{Q: How did you handle observability with Prometheus/Grafana?}\\
\textbf{A:}
\begin{itemize}[leftmargin=*]
  \item Services expose metrics endpoints via Prometheus middleware.
  \item Gateway and services map metrics endpoint and health endpoint for scraping.
  \item Docker compose includes Prometheus and Grafana for local and staging diagnostics.
  \item Logs include request context middleware for correlation and troubleshooting.
\end{itemize}

\textbf{Important snippet: metrics exposure in gateway}\newline
\begin{lstlisting}[style=code,language=C]
app.UseHttpMetrics();
app.MapGet("/health", () => "Healthy").AllowAnonymous();
app.MapMetrics("/metrics").AllowAnonymous();
\end{lstlisting}

\subsection{Failures and retries}
\textbf{Q: How do you manage failures or retries in microservices?}\\
\textbf{A:}
\begin{itemize}[leftmargin=*]
  \item We fail fast on non-transient auth/provider errors with explicit app-level codes.
  \item For external dependency calls (WSO2 token/SCIM), errors are translated into deterministic API errors.
  \item Authentication has fallback path (local JWT) during migration to reduce blast radius.
  \item Key improvement path: add policy-based retries/circuit breakers for transient network faults.
\end{itemize}

\textbf{Important snippet: failure handling for WSO2 token endpoint}\newline
\begin{lstlisting}[style=code,language=C]
try
{
    response = await _http.PostAsync(tokenEndpoint, body, ct);
}
catch (HttpRequestException ex)
{
    _logger.LogError(ex, "WSO2 token endpoint unreachable at {Endpoint}", tokenEndpoint);
    throw new AppException(503, "WSO2_UNAVAILABLE",
        "Identity server is temporarily unavailable. Please try again later.");
}
\end{lstlisting}

\section{Gateway and Security Enforcement}

\subsection{Dynamic token strategy}
\textbf{Q: How did you avoid breaking existing clients while migrating auth?}\\
\textbf{A:}
\begin{itemize}[leftmargin=*]
  \item Gateway and services use \texttt{AddPolicyScheme("DynamicJwt", ...)} to choose between local and WSO2 schemes at runtime.
  \item Token shape/issuer hints select \texttt{Wso2Jwt} for migrated tokens, else \texttt{LocalJwt}.
  \item This enabled progressive migration with low downtime risk.
\end{itemize}

\textbf{Important snippet: dynamic scheme selection}\newline
\begin{lstlisting}[style=code,language=C]
.AddPolicyScheme("DynamicJwt", "Dynamic JWT scheme", options =>
{
    options.ForwardDefaultSelector = context =>
    {
        var authHeader = context.Request.Headers.Authorization.FirstOrDefault();
        var token = ExtractBearerToken(authHeader);

        if (!string.IsNullOrWhiteSpace(token) && LooksLikeWso2Token(token, wso2ValidIssuer))
        {
            return "Wso2Jwt";
        }

        return "LocalJwt";
    };
})
\end{lstlisting}

\subsection{Rate limiting at gateway}
\textbf{Q: What protection did you add for high-frequency endpoints?}\\
\textbf{A:}
\begin{itemize}[leftmargin=*]
  \item Added fixed-window limiter for booking endpoint keyed by client IP.
  \item Configurable values: path, method, permit limit, window, rejection status code.
  \item Prevents abuse/spikes while keeping non-booking endpoints unaffected.
\end{itemize}

\textbf{Important snippet: fixed-window rate limiting}\newline
\begin{lstlisting}[style=code,language=C]
return RateLimitPartition.GetFixedWindowLimiter(partitionKey, _ =>
    new FixedWindowRateLimiterOptions
    {
        PermitLimit = bookingPermitLimit,
        Window = TimeSpan.FromMinutes(bookingWindowMinutes),
        QueueLimit = 0,
        QueueProcessingOrder = QueueProcessingOrder.OldestFirst,
        AutoReplenishment = true
    });
\end{lstlisting}

\section{Configuration and Deployment Questions}

\subsection{Environment strategy}
\textbf{Q: How do you manage environment-specific security settings?}\\
\textbf{A:}
\begin{itemize}[leftmargin=*]
  \item Keep \texttt{Wso2:Enabled=false} by default in appsettings and activate via environment variables per environment.
  \item Production requires explicit issuer/audience/client credentials and TLS enforcement.
  \item Dev profiles allow invalid certificates only when explicitly toggled.
\end{itemize}

\textbf{Important snippet: gateway security config}\newline
\begin{lstlisting}[style=code,language=]
"Wso2": {
  "Enabled": false,
  "Authority": "https://20.193.250.12:9443",
  "ValidIssuer": "https://20.193.250.12:9443/oauth2/token",
  "Audience": "REPLACE_WITH_WSO2_CLIENT_ID",
  "RequireHttpsMetadata": true,
  "AllowInvalidCertificate": false
}
\end{lstlisting}

\subsection{Migration and rollout}
\textbf{Q: How would you roll this out safely?}\\
\textbf{A:}
\begin{itemize}[leftmargin=*]
  \item Stage in non-prod with WSO2 enabled for Identity and one domain service first.
  \item Validate login, role claims, token audience/issuer, and API access policies.
  \item Gradually enable all services and monitor 401/403/error rates and latency.
  \item Keep rollback switch via \texttt{Wso2:Enabled=false} until stabilization.
\end{itemize}

\section{Quick-Fire Follow-up Questions}
\begin{enumerate}[leftmargin=*]
  \item \textbf{How do you secure SCIM admin credentials?} Use secret stores and environment injection; never commit plaintext credentials.
  \item \textbf{What if WSO2 is down?} Return deterministic service unavailable errors and alert; fallback mode can be used during migration phases.
  \item \textbf{How do you validate role claims from different issuers?} Normalize claim type mapping and policy checks for role claim variants.
  \item \textbf{How do you prevent token replay?} Short expiry, strict audience/issuer checks, TLS, and optional token introspection patterns for high-risk paths.
  \item \textbf{What is one improvement you would implement next?} Move from ROPC to Authorization Code + PKCE for stronger security posture.
\end{enumerate}

\section{Snippet Index (Repository Mapping)}
\begin{longtable}{p{0.42\linewidth} p{0.52\linewidth}}
\textbf{File} & \textbf{Why it matters} \\
\hline
\texttt{backend/QueueLanka.Identity/Services/AuthService.cs} & Login delegation to WSO2 and local fallback behavior \\
\texttt{backend/QueueLanka.Identity/Services/Wso2IdentityService.cs} & ROPC token request and SCIM provisioning implementation \\
\texttt{backend/QueueLanka.Gateway/Program.cs} & Dynamic JWT auth, JWKS authority validation, metrics, rate limiting \\
\texttt{backend/QueueLanka.Queue/Program.cs} & Service-level dynamic token scheme and claim policy handling \\
\texttt{backend/QueueLanka.Gateway/appsettings.json} & WSO2 settings and gateway policy knobs \\
\texttt{officer-api.yaml} & APIM-focused security and throttling metadata \\
\texttt{docker-compose.yml} & Service boundaries and observability stack wiring \\
\end{longtable}

\section{End-of-Interview Closing Statement}
\textbf{Suggested close:}\\
In this project, I focused on introducing standards-based identity and progressively hardening the platform without breaking existing clients. The key result is a migration-ready architecture where identity is centralized, token validation is stronger with authority/JWKS support, and governance concerns such as throttling and lifecycle can be managed at the API platform level.

\end{document}
